\documentclass[10pt,a4paper,fontset=none]{ctexart}
\usepackage{cv-template}
\setCJKmainfont{Microsoft YaHei Light}[BoldFont=Microsoft YaHei Bold,ItalicFont=Microsoft YaHei Light]
\renewcommand{\cvsmallcaps}[1]{#1}
\renewcommand{\cvnamesep}{}
\renewcommand{\cvname}{栾皓童}
\renewcommand{\cvupdated}{2026 年 9 月 15 日}
\hypersetup{pdftitle={栾皓童 - 个人简历}}

\begin{document}
\cvheader{栾}{皓童}{中国深圳}

\cvsection{研究}{概况}
自动化专业大三本科生，研究方向为欠驱动机器人手、环境自适应抓取与机器人机构设计。科研经历涵盖机械夹爪开发、机械臂视觉伺服控制与硬件原型设计，关注 VLM、VLA、WAM 在机器人可靠操作中的应用。

\cvsection{教育}{背景}
\cventry{深圳技术大学未来技术学院}{中国深圳}{自动化专业，工学学士在读（大三）}{2024.09 -- 至今}
\begin{itemize}
  \item GPA：84.17/100。研究方向：欠驱动夹爪、环境自适应机械手与机器人机构设计。
\end{itemize}
\vspace{3pt}
\cventry{清华大学 \textbar\ 深圳 X-Institute}{中国深圳}{P-SRT 项目，联合培养方向：欠驱动机器人}{2025.09 -- 至今}
\begin{itemize}
  \item 导师：张文增副教授。当前阶段：P-SRT；培养路径：X-idea $\rightarrow$ P-SRT $\rightarrow$ E-SRT $\rightarrow$ ORIC $\rightarrow$ SURF。
\end{itemize}

\cvsection{科研}{经历}
\cventry{上海人工智能实验室}{远程}{研究助理：视觉伺服控制与末端执行器设计}{2026.04 -- 至今}
\begin{itemize}
  \item 在高远研究员指导下，参与机械臂抓取的视觉伺服控制算法研究，负责臂端夹爪的机械设计。
\end{itemize}
\vspace{3pt}
\cventry{深圳 X-Institute，清华大学}{中国深圳}{P-SRT 研究者：欠驱动机器人夹爪}{2025.09 -- 至今}
\begin{itemize}
  \item 在张文增副教授指导下，独立提出双滑块自适应夹爪构想，开展机构分析与可行性研究。
  \item 设计直线运动连杆、平行四边形机构与空行程传动，实现单驱动可靠抓取；主导 DS Hand 一作设计（RAAI 2025），合作参与空行程夹爪与 GLINT 抓取提升机构研究。
\end{itemize}
\vspace{3pt}
\cventry{河套学院}{中国深圳}{量子 \& AI 冬令营}{2026.01}
\begin{itemize}
  \item 在天津大学张鹏教授指导下，研究基于 MLP/LSTM/Transformer 的量子变分线路模拟，涵盖态保真度优化、混合态模拟与物理约束训练。
\end{itemize}

\cvsection{论文}{成果}
{\fontsize{8}{10}\selectfont * 为通讯作者。在投与在审稿件单独列出。}\par
\cvsubsection{同行评审论文与已录用论文}

\cvpub{P1}{DPS Hand: A Dual-Parallelogram and Slot-Guided Gripper for Adaptive Grasping and Environmental Scooping}{\self, Haokai Ding, Tao Yang*, Wenzeng Zhang*}{IEEE CASE 2026}{已录用；RAS 三大会；中国沈阳，2026.08.17--21；第一作者。}
\cvpub{P2}{PlaceRecover: Layout-Invariant Point Cloud Recovery for Robust 3D Perception under LiDAR Placement Shifts}{Benwu Wang, Zihang Wang*, Wenkai Zhu, \self, Dong Kong, Binghao Wang, Yiming Zhou, Kailin Lyu, Teng Yan, Haoyu Wang, Yaohua Liu, Qiang Zhang}{IEEE/RSJ IROS 2026}{已录用；机器人顶会；美国匹兹堡，2026.09.27--10.01；第四作者。}
\cvpub{P3}{CSD: Conformal Selective Diagnosis for Reliable Bearing Fault Identification Under Domain Shift}{Xi Yu, Anran Lu, Keyi Chen, \self*, Yuan Gao*}{IEEE SMC 2026}{已录用；CCF C；美国贝尔维尤，2026.10.04--07；通讯作者。}
\cvpub{P4}{Domain-Semantic Subspace Stacking for Financial Fraud Detection}{\self, Minmin Chen, Zhenhui Lin, Haokai Ding, Yiwei Sheng*}{ICIC 2026}{已录用；CCF C；加拿大多伦多，2026.07.22--26；第一作者。}

\cvpub{P5}{An Underactuated Dual-Slide Adaptive Robotic Hand with Strong Environmental Adaptability}{\self, Wenzeng Zhang*, Tao Yang*}{RAAI 2025}{新加坡，2025.12.18--20；第一作者。DOI：\href{https://doi.org/10.1109/RAAI67517.2025.11423404}{10.1109/RAAI67517.2025.11423404}。}
\cvpub{P6}{An Idle-Stroke Underactuated Gripper with Independent Double-Parallelogram Linkages for Pinching and Adaptive Grasping}{Yizhe Chen, \self, Haokai Ding, Wenzeng Zhang*}{RAAI 2025}{新加坡，2025.12.18--20；第二作者。DOI：\href{https://doi.org/10.1109/RAAI67517.2025.11423397}{10.1109/RAAI67517.2025.11423397}。}
\cvpub{P7}{GLINT: An Idle-Stroke Grasp-and-Lift Hand for In-Hand Manipulation}{Haokai Ding, \self, Tao Yang*, Wenzeng Zhang*}{ICCMA 2025}{法国巴黎，2025.11.24--26；第二作者。DOI：\href{https://doi.org/10.1109/ICCMA67641.2025.11369602}{10.1109/ICCMA67641.2025.11369602}。}

\cvsubsection{在投 / 在审期刊论文}
\cvpub{M1}{Manuscript submitted to Advanced Intelligent Systems}{\self, Jie Zhang, Yuan Gao}{在投，2026}{智能系统与机器人相关方向；JCR Q1。}
\cvpub{M2}{Manuscript under review at Scientific Reports}{\self, Xi Yu, Anran Lu, Keyi Chen, Jianwu Chen}{在审，2026}{JCR Q1。}

\cvsection{在申}{专利}
\textbf{3 项发明专利}及\textbf{1 项实用新型专利}已申请，审查中。

\cvsection{荣誉}{奖励}
\cvrow{2026}{\textbf{清华大学 AIR 夏令营录取}，清华大学智能产业研究院。}
\cvrow{2026}{\textbf{CASE 2026 IEEE RAS Travel Support}（Travel Grant）。}
\cvrow{2026}{\textbf{国家级大学生创新创业训练计划项目负责人}，项目“零隙智抓”。}
\cvrow{2026}{\textbf{省级大学生创新训练计划项目成员}，项目“面向单细胞病理观测的多芯光纤动态势阱”。}
\cvrow{2026}{\textbf{浙江大学 \& 哈尔滨工业大学 SDG 全球暑期学校录取}，免学费。}
\cvrow{2026}{\textbf{量子 \& AI 冬令营入选者}，70 余名参与者中仅 9 名本科生之一，且为最年轻入选者。}
\cvrow{2024--2025}{\textbf{震雄新才院长特别奖学金}，深圳技术大学未来技术学院。}
\cvrow{2025}{\textbf{深圳技术大学科技创新“晶体奖学金”}，全校前 10 名。}
\cvrow{2024, 2025}{\textbf{深圳技术大学优秀学生干部}。}
\cvrow{2025}{\textbf{“挑战杯”中国大学生创业计划竞赛校赛优胜奖}。}
\cvrow{2025}{\textbf{“挑战杯”中国大学生课外学术科技作品竞赛校赛优胜奖}。}
\cvrow{2025}{\textbf{优秀赛会志愿者}，第十五届全国运动会暨全国第十二届残特奥会。}

\cvsection{学术}{服务与学生工作}
\cvrow{2026}{\textbf{会议审稿人}，IEEE-RAS International Conference on Humanoid Robots（Humanoids），人形机器人顶会。}
\cvrow{2026}{\textbf{会议审稿人}，IEEE/RSJ International Conference on Intelligent Robots and Systems（IROS），机器人顶会。}
\cvrow{2026}{\textbf{Workshop 审稿人}，IAB @ NeurIPS 2026 Workshop。}
\cvrow{2024--至今}{\textbf{组织部部长}，深圳技术大学未来技术学院，2024.12 起。}
\cvrow{2024--2025}{\textbf{班长兼团支书}，2024.09--2025.10。}

\cvsection{专业}{技能}
\cvskill{机构设计}{平夹耦合夹爪、直线连杆、平行四边形机构、欠驱动与腱驱动机械手；SolidWorks 建模、装配与运动仿真；Bambu Lab 3D 打印；样机迭代与抓取实验。}
\cvskill{编程与 AI}{Python；基础深度学习框架使用经验，应用于点云与量子线路模拟。}
\cvskill{学术写作}{IEEE 期刊与会议格式、投稿流程；AI 辅助文献综述。}
\cvskill{语言}{中文（母语）；英语（CET-4：580；IELTS 计划于 2026.12 考试）。}
\end{document}
